% Summary of verbilhunter2026: Omnivorous agreement, like {Uyghur} backness harmony, is a challenge for tier-based strict locality

\begin{abstract}
\textbf{Background:} Mayer and Major 2018 found that Uyghur backness harmony was not Tier-based Strictly Local, but instead Output-TSL.

\textbf{Claim:} Omnivorous Agreement in Mundari shows a similar pattern, and is also not TSL, but instead OTSL.

\textbf{Thrust:} Harmony/agreement patterns must be OTSL and not TSL if you have a hierarchy of more than two levels.

no rock paper scissors
\end{abstract}

\section*{Uyghur backness harmony}

\begin{enumerate}
  \item If there's a harmonizing vowel in the stem, the suffix should match the last one.
  \item Elif there's a harmonizing consonant in the stem, the suffix should match the last one.
  \item Else, the suffix should be a default value.
\end{enumerate}

This requires the rule, going left-to-right:

\begin{itemize}
  \item Project a harmonizing consonant only if you haven't projected a harmonizing vowel yet.
\end{itemize}

\section*{Mundari omnivorous agreement}

\begin{enumerate}
  \item If there's a first person DP as direct or indirect object, the suffix should be \textit{-iɲ}.
  \item Elif there's a second person DP as direct or indirect object, the suffix should be \textit{-m}.
  \item Else, the suffix should be \textit{-i}.
\end{enumerate}

This requires the rule, going bottom-to-top:

\begin{itemize}
  \item Project a second person DP only if you haven't projected a first person DP yet.
\end{itemize}

\section*{Definition of O-TSL}

\begin{quote}
A (\textit{j}, \textit{k})-OTSL grammar is like a \textit{k}-TSL grammar, except that the tier projection function can take into account the most recent \textit{j-1} symbols already projected to the tier when deciding whether to project a particular surface symbol.

\end{quote}

So,

A (\textit{j}, \textit{k})-OTSL language L can be defined as (Sigma, T, S, \textit{k}) where:

\begin{itemize}
  \item Sigma is a finite alphabet
  \item T is a set of tuples (sigma, tau), where sigma is in Sigma, and tau is in Union(start-symbol, Sigma) and is of length \textit{j-1}.
  \item S is a set of forbidden factors of length \textit{k}, from Sigma*
  \item Progressing from the bottom of the c-string upwards, project symbol sigma onto the tier iff there is a member of T (sigma, tau) and the tier constructed thus far ends in tau.
  \item L contains all strings in Sigma* whose tiers don't contain any factor in S
\end{itemize}

\section*{Other languages with three-level agreement hierarchies}

\begin{itemize}
  \item Hayu
  \item Onondaga
  \item Mundari (also has a Sg > Pl > Du hierarchy)
\end{itemize}

\section*{Theoretical Thrust}

Kind of backs up Graf (2022)'s Cognitive Parallelism Hypothesis:

\begin{itemize}
  \item Distinct language models (e.g. phonology, syntax) have the same formal complexity.
\end{itemize}

\section*{Thoughts about this as an MTSM}

This actually fits very well as an MTSM. Progressing bottom-up, the MTSM simply needs to read the current stored value for a given tier, and then decide whether to replace it with the current head value or not, depending on some listed hierarchy of features.
